\problemname{Treasure Map}
Paulina likes Japan, a lot.
During a vacation in Tokyo she visits an amusement park where there is a large maze.
To navigate the maze, Paulina receives a treasure map that she follows.

On the treasure map, each cell is marked with arrows to show which direction to go from that cell.

Paulina always starts in the cell in the top-left corner of the treasure map, and then follows the arrows.
In the maze there are two different goals: a piece of delicious salmon sushi, and a scary samurai.
It can also happen that the treasure map leads Paulina around in an infinite cycle of cells so she never reaches a goal.

Can you help Paulina determine which goal she reaches, or if she will walk around forever?

\section*{Input}
The first line contains an integer $R$ ($1 \leq R \leq 100$), the number of rows in the treasure map.
Then follows a line containing an integer $C$ ($1 \leq C \leq 100$), the number of columns in the treasure map.
Finally follow $R$ lines each containing $C$ characters --- the treasure map itself.

The following characters appear in the treasure map:
\begin{itemize}
\item ``\texttt{<}'' --- cell with a left arrow,
\item ``\texttt{>}'' --- cell with a right arrow,
\item ``\texttt{v}'' --- cell with a down arrow,
\item ``$\wedge$'' --- cell with an up arrow,
\item ``\texttt{A}'' --- the cell where the sushi is located,
\item ``\texttt{B}'' --- the cell where the samurai is located.
\end{itemize}

Paulina starts at the first cell in the first row of the treasure map.
The treasure map is constructed so that Paulina will never leave the maze when following the arrows.

\section*{Output}
Your program shall output a single line with the text
\begin{itemize}
\item ``\texttt{sushi}'' if she reaches the sushi by following the arrows,
\item ``\texttt{samuraj}'' if she reaches the samurai by following the arrows,
\item ``\texttt{cykel}'' if she will run around forever by following the arrows.
\end{itemize}

\section*{Scoring}
Your solution will be tested on a set of test groups.
To earn points for a group, you must pass all test cases in that group.

\noindent
\begin{tabular}{| l | l | p{12cm} |}
  \hline
  \textbf{Group} & \textbf{Points} & \textbf{Constraints} \\ \hline
  $1$    & $40$      & The treasure map is the same as on \href{https://www.progolymp.se/2018/affisch.pdf}{our poster} \\ \hline
  $2$    & $30$      & The treasure map will always lead either to the sushi or the samurai \\ \hline
  $3$    & $30$      & No additional constraints.  \\ \hline
\end{tabular}
